\chapter*{绪论}
\addcontentsline{toc}{chapter}{绪论}
\markboth{绪论}{}

本章为绪论部分。根据中国传媒大学学位论文撰写规范，章标题与第一个节标题之间必须有正文段落，不得在章标题之后直接排列节标题。本章介绍研究背景、研究问题与意义，以及研究方法与章节安排。

\section{研究背景}
本模板依据中国传媒大学研究生学位论文撰写规范制作，旨在为研究生提供符合格式要求的 \LaTeX{} 写作框架。正文采用小四号宋体，行距固定为 22 磅，左右页边距各 2.5 cm，上下页边距各 3.5 cm。论文各级标题按照“章—节—条—款”层级依次排列，格式详见后文示例。

本段为占位正文，供排版效果参考。在实际写作中，此处应填写研究背景与问题陈述，说明本研究的学术价值与实践意义。段落首行缩进两字，正文两端对齐，标点符号使用全角中文标点。

\section{研究问题与意义}
本节展示“节”级标题的排版效果。节标题为小三号黑体，居中显示，节标题上方空两行，下方空一行后接正文。以下为占位文字，实际写作时应在此阐述研究问题的界定与研究意义。

本研究拟探讨以下问题：其一，模板各级标题的层次关系如何呈现；其二，图表、引用与参考文献如何规范嵌入；其三，前置页面与正文页面的页码切换机制如何运作。以上问题将在后续章节中逐一说明。

\subsection{理论意义}
本条展示“条”级标题排版效果。条标题为四号黑体，左起空两字打印。以下为占位文字，实际写作时应在此说明本研究对相关理论的贡献与推进。

占位正文：本研究在理论层面具有以下意义。第一，补充了现有文献在格式规范方面的空白；第二，为后续研究提供了可复用的写作框架；第三，验证了 \LaTeX{} 在中文学位论文写作中的可行性。

\subsection{实践意义}
本条为第二个条级标题示例，展示同一节内多个条级标题的排版间距与视觉效果。

占位正文：本研究在实践层面具有以下价值。第一，降低了研究生在格式调整上的时间成本；第二，减少了因格式错误导致的反复修改；第三，为学院提供了统一的论文模板管理方案。

\section{研究方法与章节安排}
本节展示另一个节级标题。以下为占位正文，实际写作时应在此介绍研究方法、数据来源及各章安排。

本文共分为两章。第一章为绪论，介绍研究背景、研究问题与章节安排；第二章为主体内容，展示正文排版、图表嵌入与多级标题的综合效果。参考文献与致谢附于正文之后。

\subsubsection{文献综述方法}
本款展示“款”级标题排版效果。款标题为小四号黑体，左起空两字打印。以下为占位文字。

占位正文：本研究采用文献综述与案例分析相结合的方法。文献综述部分系统梳理了国内外关于 \LaTeX{} 学位论文模板的相关研究，案例分析部分则以本模板为对象，验证各项格式设置的实际效果。

\subsubsection{数据收集方法}
本款为第二个款级标题示例，展示同一条内多个款级标题的间距效果。

占位正文：数据收集采用问卷调查与访谈相结合的方式。问卷主要收集用户对模板易用性的评价，访谈则深入了解用户在论文写作过程中遇到的格式问题。所有数据均经过匿名化处理，仅用于本模板的改进参考。
